% !TeX root = ../初等数学笔记.tex
% 由主文档通过 \input 载入；请编译主文档。
\setcounter{chapter}{13} % 保留原书章号。
\chapter{四边形与圆}
\label{chap:em-circles}
四边形可以通过对角线分解为三角形，圆则借助半径、弦和圆周角连接不同三角形。选择辅助线时，应优先寻找已有的全等、相似和面积结构。

\section{四边形的性质与判定}
两组对边分别平行的四边形是平行四边形。作一条对角线，利用内错角与公共边可证两个三角形全等，进而得到对边相等、对角相等；比较两对角线交点附近的三角形，还可证明对角线互相平分。反之，若两对角线互相平分，由边角边全等可得对应内错角相等，从而两组对边平行。

矩形是有一个直角的平行四边形，菱形是四边相等的平行四边形，正方形同时是矩形与菱形。矩形对角线等长，菱形对角线互相垂直；结合平行四边形条件，各自的逆判定也成立。仅有对角线等长不足以保证矩形，等腰梯形就是反例。此处梯形约定为恰有一组对边平行的四边形。

\section{多边形的分割与面积}
从凸 $n$ 边形（$n\ge3$）的一个顶点向不相邻顶点连对角线，可得到 $n-2$ 个不重叠三角形，因此内角和为 $(n-2)180^\circ$。沿边界同向行走一周，转角总和为 $360^\circ$；对非凸图形须使用带方向的转角，不能把所有凹顶点都按正的小外角相加。

平行四边形面积等于底乘高，梯形两底长 $a,b$、高 $h$ 时面积为 $(a+b)h/2$，可由两份全等梯形拼成平行四边形证明。正多边形把中心与顶点相连，各小三角形具有共同的高 $\rho$，所以总面积为 $P\rho/2$，其中 $P$ 为周长，$\rho$ 为中心到边的距离。这个公式是圆面积逼近的重要桥梁。

\section{圆、弦与圆周角}
到定点 $O$ 距离等于 $r>0$ 的点集为圆，$O$ 为圆心、$r$ 为半径；圆及其内部构成圆盘。连接圆上两点的线段为弦，通过圆心的弦为直径。圆心到弦的垂线平分弦，因为垂足两侧的直角三角形斜边及一直角边相等；反向结论对非零弦也成立。

\begin{theorem}[圆周角定理]
圆上点 $P,A,B$ 两两不同，$\angle APB$ 的大小等于它所对的不含 $P$ 的弧 $AB$ 的圆心角大小的一半。
\end{theorem}
\begin{proof}
先考虑 $O$ 在角的一边上的情形，由两条半径组成的等腰三角形及外角定理得到所对圆心角为圆周角的两倍。一般情形连接 $PO$ 并延长为直径，将原角分为两个上述角或表示为它们之差；对应的弧也作同样的和或差，结论随之成立。这一处理同时覆盖圆心在角内和角外的情况。
\end{proof}
同弧上的圆周角相等，直径所对的圆周角为直角。圆内接四边形的两个对角所对弧合成整圆，故对角互补；反之，凸四边形一组对角互补可由圆周角的轨迹性质判定四点共圆。判定时须保留凸性与顶点次序。

\section{切线、割线与点的幂}
直线与圆恰有一个公共点时称为切线。切线垂直于过切点的半径：若半径不垂直，圆心到直线的垂足距离严格小于半径，直线就会有两个交点。反之，过圆上一点垂直于半径的直线也只有该交点，由勾股定理可见线上其他点到圆心更远。

设 $P$ 为圆外点，一条割线依次交圆于 $A,B$，另一条依次交于 $C,D$。由圆周角建立的相似三角形给出 $PA\cdot PB=PC\cdot PD$；若 $PT$ 为切线，则同样得到 $PT^2=PA\cdot PB$。过 $P,O$ 的割线进一步说明
\[
 PA\cdot PB=PT^2=OP^2-r^2.
\]
量 $OP^2-r^2$ 称为\emph{点的幂}（power of a point）。若 $P$ 在圆内且弦 $AB$ 经过它，无向长度则满足 $PA\cdot PB=r^2-OP^2$，其正值是点的幂的相反数，不能把圆外公式的符号直接搬入圆内。

\section{圆周长、面积与扇形}
考虑不断加倍边数的内接正多边形及对应外切正多边形。内周长增加，外周长减少，并提供共同上下界。若内接边长为 $\ell_n$、边数为 $n$、边心距为 $\rho_n$，则 $\rho_n^2=r^2-(\ell_n/2)^2$；内周长有界给出 $\ell_n\to0$，故 $\rho_n\to r$。外、内周长之比为 $r/\rho_n\to1$，因此两者趋于同一值 $C$，定义为圆周长。

相似缩放使 $C/r$ 与半径无关，记常数为 $2\pi$，得到 $C=2\pi r$。内、外面积分别为 $P_n\rho_n/2$ 与 $Q_nr/2$，都趋于 $Cr/2$；采用面积对这种穷竭逼近的相容性，圆盘面积为 $\pi r^2$。这说明公式依赖长度与面积的逼近理论，并非单凭拼接示意得出。圆心角为 $\alpha^\circ$、$0\le\alpha\le360$ 的扇形，由旋转对称、可加性及连续逼近，占整圆面积和弧长的比例均为 $\alpha/360$。

\section{例题、习题与学习检查}
\begin{example}
圆外点 $P$ 到圆心距离为 $13$，圆半径为 $5$，求切线段长；若一条割线上近交点距 $P$ 为 $8$，求远交点距离。
\end{example}
\begin{solve}
点的幂为 $13^2-5^2=144$，故切线段长为 $12$。割线关系给出 $8\cdot PB=144$，所以 $PB=18$；它大于近交点距离，符合交点次序。
\end{solve}
\begin{enumerate}
\item 菱形两对角线长为 $6,8$，求边长和面积。答：$5$ 与 $24$。
\item 半径 $3$ 的圆中，$120^\circ$ 扇形的弧长与面积是多少？答：$2\pi$ 与 $3\pi$。
\item 圆内两弦相交，其中一弦两段长为 $2,6$，另一弦一段长为 $3$，求另一段。答：$4$。
\end{enumerate}

